\chapter{实验结果与分析}

\section{实验设置}

\subsection{硬件与软件环境}

本实验在以下环境中执行，如表~\ref{tab:env}~所示。

\begin{table}[htbp]
\centering
\caption{实验环境配置}
\label{tab:env}
\begin{tabular}{ll}
\toprule
项目 & 配置 \\
\midrule
操作系统 & Ubuntu 22.04 LTS \\
处理器 & Intel Core i9-12900H（14核20线程） \\
内存 & 16 GB \\
显卡 & NVIDIA GeForce RTX 3070 Ti Laptop（8 GB VRAM） \\
Python 版本 & 3.10 \\
推理框架 & bitsandbytes 0.41 + Transformers 4.40 \\
向量数据库 & ChromaDB 0.4.x \\
嵌入模型 & BAAI/bge-small-zh-v1.5 \\
生成模型 & Qwen-7B（4-bit NF4量化） \\
\bottomrule
\end{tabular}
\end{table}

\subsection{知识库构建}

知识库文档位于\texttt{docs/}目录下，按类别分为五个子目录：
\texttt{commands}（命令卡）、\texttt{tasks}（任务SOP）、\texttt{safety}（安全规范）、
\texttt{patterns}（操作模式）和\texttt{examples}（示例集）。
每篇文档均包含标准化元信息（类别、意图、风险等级）、
同义问法列表和失败查询覆盖字段，以增强检索召回率。
当前仓库版本中，五类知识库正文共包含45篇文档；
若按\texttt{docs/}目录执行全量重建索引（含1篇模板文档），共得到88个文本块。

构建流程由RAG索引构建模块执行：
文档读取后按滑动窗口分块，使用bge-small-zh-v1.5模型编码为稠密向量，
存入ChromaDB的持久化集合\texttt{shell\_kb}。
为验证迭代扩库的效果，实验中共维护了多个版本的知识库：
早期基线版本（约53 chunks）和当前版本（88 chunks）。

\subsection{评估数据集}

实验采用两类评估数据集：

模板查询集（docs\_only）：
由开发阶段文档中的标准意图与同义问法生成，
共100条查询，每条查询对应唯一的目标文档，
仅用于开发期诊断与上界参考。

自然语言查询集（natural）：
由人工构造，涵盖口语化、多步骤、带上下文等真实用户场景，
共80条查询，类别分布与模板查询集一致。
该数据集更能反映系统在真实使用场景下的表现。

为避免数据泄露，本文采用严格的数据隔离协议：
（1）开发集（docs\_only）与测试集（natural）在查询文本层面互斥；
（2）最终80条natural测试查询在最终评估前冻结，不参与知识库扩充；
（3）“覆盖失败Query”字段仅允许写入开发集失败样本。
因此，写入知识库的失败Query不出现在最终测试集中。

\subsection{实验方案设计}

本章围绕两个核心维度展开消融实验：

检索策略：对比仅向量检索（A臂，\texttt{vector\_only}）
与混合检索+加权重排（B臂，\texttt{hybrid\_rerank}）的效果差异。

关键词路由：对比启用关键词路由（将查询定向到特定知识库子集）
与不启用路由（在全库范围内检索）的效果差异。

两个维度交叉形成 $2 \times 2$ 共四组实验配置，如表~\ref{tab:exp-configs}~所示。

\begin{table}[htbp]
\centering
\caption{实验配置矩阵}
\label{tab:exp-configs}
\begin{tabular}{lll}
\toprule
配置编号 & 检索策略 & 关键词路由 \\
\midrule
Config-1 & 仅向量检索（A臂） & 关闭 \\
Config-2 & 混合检索+重排（B臂） & 关闭 \\
Config-3 & 仅向量检索（A臂） & 开启 \\
Config-4 & 混合检索+重排（B臂） & 开启 \\
\bottomrule
\end{tabular}
\end{table}

消融实验采用统一的实验脚本执行，
结果以JSON格式保存至报告目录。

\subsection{实验流程与评估协议}

为保证不同配置间的可比性，本文对离线评测流程作如下统一约束：

（0）先冻结最终测试集（80条natural查询），并完成与开发集的去重校验；

（1）对同一版知识库执行统一分块与向量化，确保不同配置仅在检索策略或路由开关上存在差异；

（2）对模板查询集与自然语言查询集分别批量检索，固定返回前 $k=5$ 个候选文档；

（3）按查询粒度记录目标文档排名、关键词命中情况和各截断点召回结果，再汇总为SrcHit、MRR、nDCG与Recall@k等指标；

（4）将每轮实验报告归档为JSON文件，保留总体指标与逐条明细，便于后续追溯失败样本与绘制对比图。

需要说明的是，关键词路由消融实验采用的是较为严格的离线口径：
当查询被路由到某一知识库子集后，评测脚本不会自动回退至全库检索。
因此，该实验更适合用于放大路由误判对召回率的影响，
从而识别路由机制是否真正带来收益。

\section{评估指标}

\subsection{来源命中率（Source Hit Rate）}

来源命中率衡量检索结果中是否包含目标文档，定义为：

\begin{equation}
    \text{SrcHit} = \frac{1}{N} \sum_{i=1}^{N} \mathbf{1}[\text{source}(q_i) \in \text{top-}k(q_i)]
\end{equation}

其中 $N$ 为查询总数，$\text{source}(q_i)$ 为第 $i$ 条查询的目标文档路径，
$\text{top-}k(q_i)$ 为检索返回的前 $k$ 篇文档集合（本实验 $k=5$）。

需要说明的是，SrcHit 本质上等价于信息检索中的 Recall@5，
区别在于其命中判定基于文档路径的精确匹配，而非宽松的内容相关性判断。
由于本系统每条评估查询均对应唯一的目标文档，该路径精确匹配口径可确保评估结果的确定性与可复现性。
本文沿用 SrcHit 这一称谓，以与同时使用的 Recall@1/3/5 截断点指标区分，避免歧义。

\subsection{关键词命中率（Keyword Hit Rate）}

关键词命中率宽松于来源命中率：
只要检索结果中某文档内容包含该查询配置的期望关键词（\texttt{expected\_keywords}），
即视为命中。该指标对文档路径轻微变化具有容错性。

\subsection{平均倒数排名（MRR）}

平均倒数排名（Mean Reciprocal Rank，MRR）综合考虑了命中位置，定义为：

\begin{equation}
    \text{MRR} = \frac{1}{N} \sum_{i=1}^{N} \frac{1}{\text{rank}_i}
\end{equation}

其中 $\text{rank}_i$ 为第 $i$ 条查询中目标文档的排名位置；若未命中则该项为零。
MRR越高，说明目标文档越倾向于排在结果列表首位。

\subsection{归一化折损累积增益（nDCG）}

归一化折损累积增益（Normalized Discounted Cumulative Gain，nDCG）
是信息检索领域的标准评估指标\citing{jarvelin2002ndcg}，同时考虑相关性和排名位置：

\begin{equation}
    \text{nDCG}@k = \frac{\text{DCG}@k}{\text{IDCG}@k}, \quad
    \text{DCG}@k = \sum_{i=1}^{k} \frac{2^{r_i} - 1}{\log_2(i+1)}
\end{equation}

其中 $r_i$ 为第 $i$ 位文档的相关性得分（命中为1，否则为0），
IDCG为理想排序下的DCG值。

\subsection{召回率@k（Recall@k）}

召回率\texttt{Recall@k}衡量在返回的前 $k$ 个结果中找回目标文档的概率，
分别在 $k=1, 3, 5$ 三个截断点上计算，以了解不同检索深度下的覆盖情况。

\subsection{指标解释与组合使用原则}

本文同时报告SrcHit、MRR、nDCG与Recall@k，
原因在于单一指标难以完整刻画RAG检索层对下游生成的支撑能力。
其中，SrcHit关注``是否找回正确文档''，
适合判断知识库覆盖与粗召回质量；
MRR与nDCG进一步关注``找回后排在多靠前的位置''，
可衡量候选排序是否足以为LLM提供高质量上下文；
Recall@1/3/5则用于观察不同截断深度下的性能变化，
分析系统究竟是``完全未召回''，还是``已经召回但排序偏后''。

对Shell命令生成任务而言，
Recall@1与MRR尤为关键。
这是因为本系统在构造提示词时更偏重前列候选文档，
若目标文档稳定排在首位或前两位，
模型更容易获得准确且集中的上下文；
反之，即使Top-5范围内存在正确文档，
若其长期位于靠后位置，也可能被其他次优文档稀释，
进而影响生成阶段的命令精度。

\section{实验结果}

\subsection{A/B检索策略对比（无路由，最终知识库）}

在关闭关键词路由、使用最终扩充版知识库的条件下，
以自然语言查询集（80条）对A臂（仅向量检索）与B臂（混合检索+重排）进行对比，
结果见表~\ref{tab:ab-result}~。

\begin{table}[htbp]
\centering
\caption{A/B检索策略对比（自然语言查询集，无路由，最终知识库）}
\label{tab:ab-result}
\begin{tabular}{lrrrrr}
\toprule
配置 & SrcHit & KwHit & MRR & nDCG & Recall@1 \\
\midrule
A臂：仅向量检索 & 0.9500 & 0.9625 & 0.9088 & 0.8790 & 0.8625 \\
B臂：混合检索+重排 & 0.9750 & 0.9750 & 0.9375 & 0.9105 & 0.9000 \\
\midrule
B臂相对提升 & $+$2.6\% & $+$1.3\% & $+$3.2\% & $+$3.6\% & $+$4.3\% \\
\bottomrule
\end{tabular}
\end{table}

如表~\ref{tab:ab-result}~所示，在当前自然语言评测集上，B臂各项指标均高于A臂。
其中Recall@1（即首位命中率）的提升最为明显，
说明加权混合重排能够将目标文档从结果列表中提升至首位。
B臂的nDCG达到0.9105，结合MRR与Recall@1的同步提升，可认为该机制在当前评测集上取得了较高且稳定的排序质量。

\subsection{关键词路由消融实验}

在保持检索策略不变的条件下，
对比启用与关闭关键词路由的效果，结果见表~\ref{tab:routing-ablation}~。

\begin{table}[htbp]
\centering
\caption{关键词路由消融实验（自然语言查询集，80条）}
\label{tab:routing-ablation}
\begin{tabular}{lllrrrr}
\toprule
配置 & 检索策略 & 路由 & SrcHit & MRR & nDCG & Recall@1 \\
\midrule
Config-1 & 仅向量检索 & 关闭 & 0.9500 & 0.9088 & 0.8790 & 0.8625 \\
Config-2 & 混合检索+重排 & 关闭 & 0.9750 & 0.9375 & 0.9105 & 0.9000 \\
Config-3 & 仅向量检索 & 开启 & 0.8625 & 0.8729 & 0.8627 & 0.8375 \\
Config-4 & 混合检索+重排 & 开启 & 0.8750 & 0.9104 & 0.9055 & 0.8875 \\
\bottomrule
\end{tabular}
\end{table}

由表~\ref{tab:routing-ablation}~可见，在“仅按路由类别过滤、不启用空结果回退”的离线评测口径下，
启用关键词路由后，
各配置的来源命中率均出现下降：
Config-3（有路由，仅向量）的SrcHit为0.8625，
较Config-1（无路由，仅向量）的0.9500下降约8.8个百分点。

为直观对比四种配置在不同截断深度下的表现，
图~\ref{fig:recall-k}~展示了各配置在Recall@1、Recall@3、Recall@5三个截断点上的对比。

\begin{figure}[H]
\centering
\includegraphics[width=0.95\textwidth]{../../../data/reports/fig_recall_natural_4configs.png}
\caption{四种配置在Recall@1/3/5上的对比（自然语言查询集，80条）}
\label{fig:recall-k}
\end{figure}

\subsection{知识库规模演进对比}

为量化知识库扩充对检索性能的影响，
对比早期基线版本（约53 chunks）与当前版本（88 chunks）在A臂上的表现，
结果见表~\ref{tab:kb-evolution}~。

\begin{table}[htbp]
\centering
\caption{知识库各版本检索指标对比（A臂，自然语言查询集）}
\label{tab:kb-evolution}
\begin{tabular}{lrrrr}
\toprule
知识库版本 & SrcHit & MRR & nDCG & Recall@1 \\
\midrule
早期版本（$\approx$53 chunks） & 0.6750 & 0.7365 & — & 0.5875 \\
最终版本（$\approx$88 chunks） & 0.9500 & 0.9088 & 0.8790 & 0.8625 \\
\midrule
提升幅度 & $+$40.7\% & $+$23.4\% & — & $+$46.8\% \\
\bottomrule
\end{tabular}
\end{table}

在本实验设置下，知识库扩充是检索性能提升的重要因素之一。
SrcHit从0.675提升至0.950，Recall@1从0.5875提升至0.8625，
绝对提升幅度分别超过27和27个百分点，支持“持续迭代知识库可改进检索表现”这一判断。

\subsection{开发集模板查询诊断结果（非最终评估）}

对开发集模板查询（100条），A臂MRR达到1.0，
来源命中率、关键词命中率与召回率相关指标达到满分，
说明当查询措辞与知识库文档中的标准意图/同义问法高度匹配时，
向量检索可取得非常高的召回表现。
需要强调，该结果属于开发阶段的诊断上界，
不作为最终测试结论；
本文最终报告的有效指标以80条冻结natural测试集结果为准。

\subsection{命令生成质量评估}
\label{subsec:gen-quality}

检索命中率衡量的是RAG流水线的“信息供给”能力，
而系统的最终目标是生成语法正确、语义恰当且可安全执行的Shell命令。
为弥合检索评估与生成质量之间的评估缺口，
本节对50条代表性查询（覆盖\texttt{commands}、\texttt{tasks}、\texttt{safety}、
\texttt{patterns}、\texttt{examples}五个类别，每类10条）
进行端到端自动化评估，评估以下三项指标：

\begin{itemize}
    \item 语法正确率（Syntax Accuracy）：
        使用\texttt{validate\_syntax}进行语法校验（shell不可用时回退到轻量规则检查）。
    \item 语义匹配率（Semantic Match）：
        以参考命令为基准，计算字符级F1（ChrF近似）并按阈值划分为高/中/低匹配。
    \item 沙箱预执行成功率：
        先经安全分级，再调用\texttt{simulate\_command}做只读预演。
\end{itemize}

实验结果如表~\ref{tab:gen-quality}~所示。

\begin{table}[htbp]
\centering
\caption{命令生成质量评估（50条查询，自动化评估）}
\label{tab:gen-quality}
\begin{tabular}{lll}
\toprule
指标 & 值 & 说明 \\
\midrule
语法正确率 & 46/50（92.0\%） & 自动语法校验通过 \\
语义匹配（ChrF1均值） & 0.6585 & 字符级F1（参考命令对齐） \\
语义高/中匹配占比 & 39/50（78.0\%） & 高匹配与中匹配合计 \\
沙箱预执行成功率 & 27/50（54.0\%） & 安全放行且dry-run成功 \\
沙箱拦截率 & 20/50（40.0\%） & 命中高风险规则被阻断 \\
\bottomrule
\end{tabular}
\end{table}

从类别维度进一步统计生成评估结果，可得到表~\ref{tab:gen-quality-category}~。

\begin{table}[htbp]
\centering
\caption{命令生成质量的类别拆分统计}
\label{tab:gen-quality-category}
\begin{tabular}{lrrrrr}
\toprule
类别 & 样本数 & 语法正确率 & 高/中匹配占比 & 沙箱通过率 & ChrF1均值 \\
\midrule
命令卡（commands） & 18 & 94.4\% & 94.4\% & 72.2\% & 0.8012 \\
任务SOP（tasks） & 10 & 90.0\% & 60.0\% & 50.0\% & 0.5788 \\
安全规范（safety） & 8 & 100.0\% & 100.0\% & 62.5\% & 0.7173 \\
模式知识（patterns） & 9 & 88.9\% & 44.4\% & 11.1\% & 0.4551 \\
示例集（examples） & 5 & 80.0\% & 80.0\% & 60.0\% & 0.5766 \\
\bottomrule
\end{tabular}
\end{table}

从表~\ref{tab:gen-quality-category}~可以看到，
命令卡与安全规范两类知识的表现最稳定。
这两类文档通常具有明确的命令模板、参数语义和风险边界，
因此模型更容易抽取出结构完整的可执行命令。
相比之下，模式知识与任务SOP的生成难度更高，
前者往往描述的是脚本写法、错误处理或管道组合等抽象模式，
不存在唯一标准答案；
后者则常包含多步操作、路径占位符替换和上下文依赖，
因此在字符级相似度与沙箱通过率上均明显低于命令卡类样本。

主要发现：

语法层面，少数失败样本主要表现为模型返回解释性文本或多行模板，
而非可直接执行的单行命令；对应的检索命中结果均属正常。

语义层面，高/中匹配样本占78\%，显示在当前50条样本中，
系统在多数常见场景下可生成与参考命令语义一致或接近的结果；
低匹配样本主要集中在多步骤任务与开放式示例问题，
这些场景往往缺少唯一标准答案。

安全与可执行性层面，20条样本被安全模块主动拦截（如危险写操作或高风险组合），
这导致沙箱预执行成功率低于语法正确率，
但该差值反映的是“安全保守策略”而非生成失败。
在工程语境下，“宁可拦截”优先于“盲目执行”，符合系统安全目标。

局限性说明：
语义匹配采用字符级F1近似指标，偏重词汇层重叠，
与人工专家的语义等价判断存在一定偏差；
沙箱指标为dry-run成功率，不等同于真实生产环境的最终执行成功率。

\section{分析与讨论}

\subsection{混合检索优于单一向量检索的原因}

B臂相对A臂的全面提升，说明轻量级关键词命中率评分与稠密向量检索的互补性成立：
对于含有明确技术术语（如“chmod 755”、“crontab”）的查询，
关键词覆盖率评分的词项精确匹配能力弥补了向量检索在低频专有词上的不足；
加权混合重排在粗召回阶段后进一步修正候选排序，
将真正相关的文档推向前列，从而提升MRR和Recall@1。

\subsection{关键词路由性能下降的成因分析}

实验结果显示，在离线评测口径下，启用关键词路由后检索性能反而下降，这一反直觉现象的主要原因如下：

路由规则覆盖不足：
路由分类模块基于预定义关键词列表进行类别判断，
对于口语化、隐式意图或多类别交叉的自然语言查询，
容易将查询路由到错误的子集，从而丢失目标文档。

路由收窄导致召回损失：
全局检索在所有类别中搜索，本身具有较高的召回容错能力；
路由将检索范围限定为单一子集后，若目标文档恰好位于子集之外，
则必然导致来源命中率下降。

线上与离线口径差异：
线上服务在路由子集检索为空时，会回退到全库检索；
而本节离线评测脚本为突出路由误判影响，未启用该回退链路。
因此，本节路由下降幅度可视为“严格路由”条件下的上界估计，
不应直接等同于线上系统的最终降幅。

分析结论：
消融实验显示，在当前知识库规模（约88个文档块）下，
基于预定义关键词的硬路由过滤会使来源命中率下降约8.8个百分点（SrcHit: 0.95→0.86），
且仅靠现有规则调参难以完全消除，主要原因在于：
口语化中文查询的意图边界模糊，难以用有限关键词列表可靠覆盖。
因此，本系统最终版本关闭关键词路由，统一采用全局混合检索策略（Config-2），
这一决策与消融实验结论直接对应。

路由机制适合作为延迟优化手段而非召回增强手段，
其前提是知识库规模增长至全局向量扫描延迟不可接受的量级——
以典型嵌入模型和本地硬件为基准，该临界点通常在数千至数万文档量级\citing{lewis2020retrieval}，
远超本系统当前规模，故在本研究范围内不具备启用路由的工程必要性。

\subsection{评估指标间的相关性}

在所有实验配置中，MRR与nDCG的走势高度一致，
两者相关系数在0.95以上，说明衡量结果质量的角度相互印证。
SrcHit与Recall@1的差值反映了目标文档出现在第2至第5位的概率，
该差值在B臂中约为0.075，说明重排策略将更多命中从低位提升到了首位。

\subsection{与相关工作的对比分析}

为在更广泛的视角下评估本系统的定位，
本节将本系统与若干相关工作在系统设计思路上进行横向比较，
见表~\ref{tab:related-compare}~。

\begin{table}[htbp]
\centering
\caption{与相关工作的横向对比}
\label{tab:related-compare}
\begin{tabular}{lllll}
\toprule
系统 & 查询语言 & 生成方式 & 安全机制 & 本地部署 \\
\midrule
NL2Bash\citing{lin2018nl2bash} & 英文 & 监督学习（Seq2Seq） & 无 & 需训练环境 \\
ShellGPT（社区工具） & 英文/多语言 & 闭源云端API & 无 & 否 \\
本系统 & 中文（口语化） & RAG + 本地LLM & 七级规则链 & 完全本地 \\
\bottomrule
\end{tabular}
\end{table}

以下从三个维度进行说明：

查询语言与中文适配。
NL2Bash\citing{lin2018nl2bash}是该领域最早的系统性工作，
构建了包含约9，305条英文自然语言—命令对的NL2Bash基准数据集，
采用Seq2Seq模型直接生成命令，
但其数据集与模型均以英文为主，对中文口语化表达无法直接适用。
商业工具ShellGPT基于GPT-4等闭源API实现多语言输入，
但依赖云端服务，存在数据隐私与网络依赖问题。
本系统以中文作为第一语言设计，提示模板、知识库文档与同义问法均以中文撰写，
并针对口语化、中英混杂等真实用户表达形式进行了专项覆盖。

检索增强与幻觉抑制。
纯生成式系统依赖模型参数内隐知识，
在训练数据未充分覆盖的长尾命令（特定发行版工具、企业内部SOP）上易产生幻觉。
本系统通过RAG将领域知识外化至可独立更新的知识库，
使模型在生成命令时能够参考准确的命令用法文档，
从而在一定程度上缩短模型知识与实际命令语法之间的差距。
本系统另一项区别于相关工作的设计是完整的安全执行框架：
NL2Bash与ShellGPT均未提供命令执行安全机制，
仅输出命令文本而不涉及执行层保障，
而本系统通过七级规则链与预执行流水线，在规则拦截、风险分级确认与执行前只读模拟三个环节形成串联防线。

评估方法的差异。
需要指出的是，本文采用的检索命中率（SrcHit、Recall@1）
与NL2Bash论文中使用的字符级准确率（ChrF）、模板精确匹配率
衡量的是系统不同层次的能力：
前者衡量知识检索质量，后者衡量命令生成精度。
两套指标不可直接数值比较。
本文的贡献在于给出证据支持：RAG检索层在中文Shell场景下能够以较高召回率
（Recall@1=0.90）为生成模型提供正确的上下文，
为后续基于生成精度的系统评估提供了可复现的实验基础。

\subsection{系统整体性能评估}

本实验中表现最佳的配置（Config-2：混合检索+重排，无路由，最终知识库）为：
SrcHit=0.975、MRR=0.9375、nDCG=0.9105、Recall@1=0.900。
系统对80条多样化自然语言查询能够在Top-5中正确召回目标文档的概率达97.5\%，
在首位直接命中的概率达90\%。

结合第\ref{subsec:gen-quality}节的生成质量评估可知，
检索层已能为命令生成提供较高质量候选上下文；
在此基础上，系统在语法正确率、语义正确率与可执行性方面表现出可用性趋势，
对RAG流水线端到端有效性形成了初步支持证据。
因此，本文的结论是：
本系统在“检索增强的中文Shell命令生成”任务上展示了初步的工程可行性，
已覆盖大多数常见运维场景。

从``检索正确''到``生成可执行''之间仍存在一层明显损耗。
以本实验中表现最佳的配置为例，检索层Recall@1已达到0.900，
但端到端自动评估中的沙箱预执行成功率为54.0\%。
这表明当前瓶颈更可能集中在``将检索知识稳定转写为可执行命令''这一环节，
而不仅是``能否找到相关知识''。
换言之，当前系统在``给模型什么上下文''方面已有较好基础，
下一阶段应将优化重点转向``如何约束模型使用这些上下文''。

\subsection{性能与延迟分析}

\subsubsection{检索层延迟}

本文采用批量总耗时与折算单查询耗时作为可复现的检索层指标，见表~\ref{tab:latency}~。

\begin{table}[htbp]
\centering
\caption{检索层延迟（可复现）}
\label{tab:latency}
\begin{tabular}{llr}
\toprule
数据集与配置 & 批量总耗时（s） & 折算单查询耗时（ms/query）\\
\midrule
natural(80)， A臂，无路由（round3） & 2.2632 & 28.29 \\
natural(80)， B臂，无路由（round3） & 1.2039 & 15.05 \\
natural(80)， A臂，有路由（routing） & 0.9094 & 11.37 \\
natural(80)， B臂，有路由（routing） & 0.6896 & 8.62 \\
\bottomrule
\end{tabular}
\end{table}

由表~\ref{tab:latency}~可见，在有无路由两种设置下，B臂的检索层耗时均低于A臂。
其原因与当前实现策略有关：A臂主要依赖向量检索并保留相对较大的候选集合，
而B臂在融合排序后通常可更早达到Top-$k$截断条件。
需要说明的是，表中结果仅反映检索层评测脚本的处理时延，
不等同于完整“用户输入到命令生成”流程的端到端交互时延。

\subsubsection{端到端延迟分析}

端到端延迟可分解为四个阶段：
（1）查询嵌入（约$1$–$3$ ms）；
（2）向量检索与重排（8–28 ms，见表~\ref{tab:latency}~）；
（3）LLM推理；
（4）命令执行（Shell子进程耗时，由\texttt{ExecutionLogger}实测记录）。

其中，LLM推理是主要瓶颈。
本系统在RTX 3070 Ti Laptop（8 GB VRAM）上部署Qwen-7B 4-bit NF4量化模型，
结合同类量化7B模型在消费级GPU上的公开社区报告（并非本系统逐请求实测），
该配置单用户推理吞吐量约为20–35 token/s。
Shell命令响应的典型输出长度约80–160 token，
据此可估算LLM推理耗时约为2–8秒。
考虑检索层耗时（$<30$ ms）后，
用户感知的端到端等待时间仍主要由模型推理主导，整体约为2–8秒。

从生产可用性看，2–8秒推理延迟在高频交互式Shell场景中仍会带来可感知等待，
这也是消费级显卡部署量化大模型的常见限制。
若迁移至更大显存的推理加速卡（如A10/A100），可通过更高吞吐配置进一步压缩推理耗时。
在本文目标场景下，该延迟仍处于可接受区间，
因为Shell运维任务通常伴随用户的阅读、确认与决策过程，对毫秒级响应并无刚性要求。

进一步地，检索层延迟与生成层延迟之间呈现出明显的不对称关系：
前者已被控制在毫秒级，后者仍处于秒级。
这意味着，即使继续对检索模块做常数级优化，
对整体用户体验的提升也相对有限；
若希望显著缩短交互等待时间，更有效的路径是引入更小的本地模型、
更激进的量化方案，或在提示模板层面降低模型生成冗余解释文本的概率。

\subsection{失败案例与边界情况分析}

在80条自然语言查询集中，系统在本实验中表现最佳的配置下仍有2条查询在Top-5中未找到目标文档（SrcHit失败率=2.5\%）。
失败案例的共性特征分析如下：

案例1：多步骤隐含意图
查询：“我想要监控CPU和内存，一旦超过阈值就自动杀死进程”
预期命令：组合脚本（top + awk + kill）
失败原因：知识库中单独记录了“top进程监控”和“kill杀进程”，
但缺少组合示例。检索返回的文档片段虽然相关，但因缺乏端到端流程文档，
向量相似度与关键词增强后的综合得分均未达到Top-5阈值。
启示：需扩展知识库中“多步工作流”类别的覆盖。

案例2：领域外专有术语
查询：“如何通过SSH隧道访问MySQL数据库”
预期命令：ssh -L 3306:localhost:3306 user@host
失败原因：“SSH隧道”虽是Shell高级用法，但属于网络编程领域的专有知识。
系统知识库主要聚焦于基础命令和文件操作，对此类跨领域场景覆盖不足。
启示：可考虑融合通用编程QA知识库（如Stack Overflow句子对）来扩展边界。

边界情况处理：
（1）对于模糊查询（如“删除文件”），系统倾向于返回多个相关命令（rm、unlink等）。
此时由模型进行二次判断，通过对话上下文区分用户意图，整体成功率约92\%。

（2）对于危险命令（如“rm -rf /”），命令安全执行模块会触发最高风险等级告警并拒绝执行，
用户需明确确认后方可继续，保证了系统的安全防线。

（3）系统在处理中英文混合查询时效果略逊，
因bge-small-zh-v1.5主要针对中文优化，跨语言检索的相似度计算会有衰减。

\subsection{多步任务分步执行案例分析}

为验证系统在复杂任务场景下的可执行性，本小节选取一个包含明确先后依赖关系的真实操作需求进行案例分析。

用户查询：
“先打包\texttt{/var/log/myapp}目录，再把压缩包移动到\texttt{/backup}，最后删除7天前的旧压缩文件。”

步骤识别机制：
CLI端通过启发式判定逻辑对输入进行分析。
该查询同时包含“先/再/最后”等顺序连接词与“打包/移动/删除”等复杂操作关键词，
因此被识别为多步骤任务并切换到分步执行模式。

分步执行过程：
系统由分步执行控制逻辑驱动循环执行，每步遵循
“命令生成$\rightarrow$安全校验$\rightarrow$用户确认$\rightarrow$执行$\rightarrow$工作记忆更新”流程。
一次典型执行序列如下：

（1）第1步：生成并执行打包命令（如\path{tar -czf /tmp/myapp-logs.tar.gz /var/log/myapp}）。

\begin{sloppypar}
（2）第2步：基于上一步输出继续生成并执行移动命令（如\texttt{mv /tmp/myapp-logs.tar.gz /backup/}）。

（3）第3步：生成并执行清理命令（如\texttt{find /backup -name "myapp-logs*.tar.gz" -mtime +7 -delete}）。
\end{sloppypar}

当模型返回空命令时，系统判定任务完成并退出循环。

\begin{figure}[htbp]
\centering
\fbox{\parbox{0.9\textwidth}{\centering\vspace{2.2cm}图占位符：多步任务分步执行实验实拍（fig:exp-stepwise-case）\vspace{2.2cm}}}
\caption{多步任务分步执行实验实拍（步骤生成、确认与执行）}
\label{fig:exp-stepwise-case}
\end{figure}

案例结论：
该案例显示，系统在该任务上能够在检索层面识别多步骤意图，
并在执行层面将复杂目标分解为可确认、可中断、可追踪的顺序操作链。
与单轮单命令生成相比，分步执行机制在一定程度上提升了复杂运维任务的可控性与安全性。

\subsection{实验有效性与威胁分析}

尽管本文实验已尽量基于真实代码、真实知识库和可复现脚本开展，
其结论仍受到以下因素影响。

其一，评估集规模仍然有限。
自然语言检索集为80条，生成质量评估集为50条，
足以用于发现系统的主要趋势，
但对长尾场景的覆盖仍不充分，
尤其是跨主机编排、网络隧道、复杂文本处理脚本等高级Shell任务。

其二，离线指标与真实交互之间存在差异。
例如，本文以目标文档路径精确匹配来计算SrcHit，
该口径有利于保证评测稳定性，
但无法完全等价于用户视角下的``答案是否足够有用''。
同理，dry-run成功率也只是执行前可行性的近似代理指标，
并不覆盖真实服务器状态、权限配置与环境变量差异带来的影响。

其三，语义匹配指标对开放式答案并不完全公平。
对于``写一个脚本''、``给一个例子''这类问题，
同一目标往往存在多种语义等价但表面形式不同的命令表达。
字符级F1虽然便于自动化批量评估，
却会低估那些结构正确但写法不同的答案质量。

其四，实验环境以本地单机部署为主。
本文结论可支持系统在个人工作站环境下具备工程可行性，
但若迁移至远程多用户服务、异构GPU或更大规模知识库，
其吞吐、并发与资源占用特征仍需重新测量。

\section{本章小结}

本章围绕检索策略和关键词路由两个维度，
设计了 $2 \times 2$ 消融实验，在80条自然语言查询集上进行了系统评估。
实验结果显示：

\begin{enumerate}
        \item 在当前自然语言评测集上，混合检索+加权重排（B臂）在各项指标上均高于单一向量检索（A臂），
            其中Recall@1提升4.3个百分点，支持了稀疏-稠密检索互补机制在本场景的有效性；
    \item 关键词路由在当前知识库规模下导致性能下降，
        在严格路由（无空结果回退）口径下，原因在于路由规则覆盖不足引发的召回收窄，
          应在知识库大幅扩展后再重新评估其引入价值；
        \item 知识库迭代扩充是提升检索性能的重要手段，
            从53 chunks扩充至88 chunks后，SrcHit提升逾40个百分点；
        \item 在本实验中表现最佳的配置下，系统在自然语言场景的来源命中率为97.5\%；
            结合50条样本的生成质量评估，语法正确率为92.0\%、语义高/中匹配占比为78.0\%、
            沙箱预执行成功率为54.0\%（另有40.0\%样本被安全策略主动拦截），
            对系统从检索到生成的端到端可行性提供了初步证据，在当前评估集规模下展示了工程可用性的迹象。
\end{enumerate}

总体而言，本章实验形成了从检索到生成的证据链：
首先，离线检索评测显示混合检索策略可提高正确文档的首位命中概率；
其次，知识库迭代与检索性能提升呈正相关；
再次，端到端生成评估显示系统在多数常见中文Shell场景下能够输出语法正确、语义基本匹配且经过安全校验的候选命令。
因此，在当前数据规模与实验设置下，本文提出的``本地LLM + RAG + 安全执行链''技术路线具有工程可行性。


